\documentclass{beamer}

% Standard Packages
\usepackage[utf8]{inputenc}
\usepackage[T1]{fontenc}

\usefonttheme{serif}

\usepackage{hyperref}
\hypersetup{
	pdftitle={Slides Template},
	pdfauthor={Euan J. Mendoza},
	colorlinks=true,
	linkcolor=black,
	citecolor=blue,
	pdfstartview=FitH,
	pdfpagemode=UseOutlines,
	pdfpagelayout=OneColumn,
}

% Personal Macros
\usepackage{em-mathtools}
% Note: em-mathtools detects beamer and suppresses its own amsthm environments.
% Use beamer's built-in theorem/block environments instead.

% \usepackage{todonotes}   % todonotes conflicts with beamer; uncomment with care.

% Bibliography Package
% \usepackage[style=alphabetic,maxbibnames=6]{biblatex}
% \addbibresource{references.bib}

% \usepackage[capitalise]{cleveref}

% Tikz setup
% \usepackage{tikz}
% \usetikzlibrary{cd}

\title{Slides Template}
\subtitle{A Beamer Example}
\author{Euan Mendoza}
\date{\today}

\begin{document}

% ---- Title frame ----
\begin{frame}
    \titlepage
\end{frame}

% ---- Table of contents ----
\begin{frame}{Outline}
    \tableofcontents
\end{frame}

%% ============================================================
\section{Beamer Basics}
%% ============================================================

\begin{frame}{Frames and Pauses}
    Beamer slides are built from \texttt{frame} environments.
    Use \texttt{\textbackslash pause} to reveal content step by step:

    \begin{itemize}
        \item First point.
        \pause
        \item Second point, revealed on click.
        \pause
        \item Third point, revealed on the next click.
    \end{itemize}
\end{frame}

\begin{frame}{Overlay Specifications}
    Finer control with overlay specs:

    \begin{itemize}
        \item<1-> Visible from slide 1 onward.
        \item<2-> Visible from slide 2 onward.
        \item<3> Visible on slide 3 only.
    \end{itemize}

    \vspace{1em}
    \only<1>{This text appears on slide 1 only.}
    \only<2->{This text appears from slide 2 onward.}
\end{frame}

\begin{frame}{Columns}
    \begin{columns}[T]
        \begin{column}{0.48\textwidth}
            \textbf{Left column} \\[0.5em]
            Use columns to place content side by side, e.g.\ a statement
            on the left and a diagram on the right.
        \end{column}
        \begin{column}{0.48\textwidth}
            \textbf{Right column} \\[0.5em]
            \[
                \R^{n} = V \oplus V^{\perp}
            \]
        \end{column}
    \end{columns}
\end{frame}

%% ============================================================
\section{Theorem Environments}
%% ============================================================

% Beamer provides its own theorem-style blocks.
% em-mathtools suppresses its amsthm definitions automatically.

\begin{frame}{Built-in Beamer Blocks}
    \begin{theorem}[Cayley]
        Every group $G$ of order $n$ embeds into $\Sym(\{1,\ldots,n\})$.
    \end{theorem}

    \pause

    \begin{proof}
        The left-regular action $g \mapsto (x \mapsto gx)$ gives an
        injective homomorphism $G \hookrightarrow \Sym(G)$.
    \end{proof}
\end{frame}

\begin{frame}{Definition and Example Blocks}
    \begin{definition}
        The \emph{Grassmannian} $\Gr(k, \F^{n})$ is the set of all
        $k$-dimensional subspaces of $\F^{n}$.
    \end{definition}

    \pause

    \begin{example}
        $\Gr(1, \R^{3})$ is the real projective plane $\mathbb{P}^{2}(\R)$.
    \end{example}
\end{frame}

\begin{frame}{Alert and Example Blocks}
    \begin{block}{Standard block}
        Use for remarks, notes, and neutral statements.
    \end{block}

    \begin{alertblock}{Alert block}
        Use to highlight warnings, important caveats, or open problems.
    \end{alertblock}

    \begin{exampleblock}{Example block}
        Use for worked examples and concrete illustrations.
    \end{exampleblock}
\end{frame}

%% ============================================================
\section{em-mathtools Macros}
%% ============================================================

\begin{frame}{Number Sets and Fields}
    Standard blackboard-bold sets:
    \[
        \N \subset \Z \subset \Q \subset \R \subset \C.
    \]
    Generic field $\F$ or $\K$; e.g.\ $\F_{q}^{n}$ is an $n$-dimensional
    vector space over the finite field with $q$ elements.
\end{frame}

\begin{frame}{Linear Algebra}
    \begin{itemize}
        \item Inner product: $\ip{u}{v}$
        \item Transpose: $\tr{A}$, \quad Trace: $\Tr(A)$
        \item Rank: $\rk(A)$
        \item Span: $\linspan\{v_{1}, \ldots, v_{k}\}$
        \item Vectors: $\vec{v}$, $\vzero$, $\vA$, $\vB$
        \item Calligraphic spaces: $\cH$, $\cA$, $\cB$
    \end{itemize}

    \pause

    \begin{block}{Rank-Nullity}
        For $T : \F^{n} \to \F^{m}$,
        \[
            \rk(T) + \dim\ker(T) = n.
        \]
    \end{block}
\end{frame}

\begin{frame}{Group Theory}
    \begin{itemize}
        \item Symmetric group: $\Sym(X)$
        \item Automorphism group: $\Aut(X)$
        \item General linear group: $\GL(n, \F)$
        \item Special orthogonal group: $\SO(n)$
        \item Unitary group: $\U(n)$
    \end{itemize}
\end{frame}

\begin{frame}{Quantum Mechanics}
    Dirac notation:
    \[
        \braket{\psi}{\phi} = \bra{\psi}\ket{\phi}
        \qquad
        \ketbra{\psi}{\psi} \text{ (projector)}
    \]

    The Hamiltonian $H$ acts on states $\ket{\psi} \in \cH$:
    \[
        H\ket{\psi} = E\ket{\psi}.
    \]
\end{frame}

\begin{frame}{Probability and Complexity}
    \begin{itemize}
        \item Expectation: $\E[X]$, \quad Variance: $\Var[X]$
        \item Polynomial growth: $\poly(n)$, $\polylog(n)$
    \end{itemize}

    \pause

    \begin{alertblock}{Key Open Problem}
        Graph Isomorphism (\GIp) is in \NP\ and in \AM\ $\cap$ \SZK,
        but is not known to be \NPc.
    \end{alertblock}
\end{frame}

%% ---- Final frame ----
\begin{frame}
    \centering
    \Large Thank you.
\end{frame}

\end{document}
